\documentclass[14pt,a4paper]{article}


\usepackage[utf8]{inputenc}
\usepackage[T2A]{fontenc}
\usepackage[russian]{babel}


\usepackage{times}   
\usepackage[a4paper,top=20mm,bottom=20mm,left=30mm,right=15mm]{geometry}
\usepackage{setspace}
% для pdfLaTeX
%\usepackage{times}  % стандартный Times, не mathptmx
%\usepackage{tempora}

\usepackage[utf8]{inputenc}
\usepackage[T2A]{fontenc}
\usepackage[russian]{babel}
\usepackage{graphicx}
\usepackage{amsmath,amssymb}
\usepackage{titlesec}
\usepackage{indentfirst} % Для отступов в первом абзаце
\usepackage{tocloft}


% Шрифт и отступ в содержании
\renewcommand{\cftsecfont}{\bfseries} % жирный шрифт для секций
\renewcommand{\cftsecpagefont}{\bfseries} % жирный номер страницы
\renewcommand{\cftsecleader}{\cftdotfill{\cftdotsep}} % точки
% Отступы
\setlength{\cftbeforesecskip}{0.2em}

% Выравнивание и заголовок
\renewcommand{\contentsname}{Содержание}
\renewcommand{\cftsecaftersnum}{\hspace{1em}} % пробел после номера


% Настройки
\onehalfspacing
\setlength{\parindent}{1.25cm} % Абзацный отступ 1.25 см

% Настройка стилей заголовков
\titleformat{\section}
  {\normalfont\fontsize{14}{16}\bfseries\centering} % стиль + выравнивание по центру
  {\thesection}{1em}{}


\titleformat{\subsection}{\normalfont\fontsize{14}{16}\bfseries\centering}{\thesubsection}{1em}{}
\titlespacing*{\section}{0pt}{3ex plus 1ex minus .2ex}{2ex plus .2ex}
\titleformat{\subsubsection}{\normalfont\fontsize{14}{16}\bfseries\centering}{\thesubsubsection}{1em}{}

% Перенумерация подзаголовков в формате Глава.Параграф
\renewcommand{\thesubsection}{\thesection.\arabic{subsection}}

% Настройка нумерации страниц
\usepackage{fancyhdr}
\pagestyle{fancy}
\fancyhf{}
\cfoot{\thepage}
\renewcommand{\headrulewidth}{0pt}

% Настройка подписей к рисункам и таблицам
\usepackage{caption}
\captionsetup[figure]{justification=centering,labelsep=endash}
\captionsetup[table]{justification=centering,labelsep=endash}



% NEW-----------------------------------------
\usepackage{amsmath}   % расширенные возможности для формул
\usepackage{amsthm}    % окружения theorem, proof и т.п.
\usepackage{amssymb}   % дополнительные математические символы (\mathbb, \varnothing и т.д.)
\usepackage{float}

  
\input{preamble.tex}
%%% Титульный лист
% \input{title}
% \setcounter{page}{2}



\renewcommand{\abstractname}{Аннотация}

\title{Прогнозирование социальных тенденций с высокой волатильностью}

% \author{
%     Мария Никитина \\
% 	\texttt{nikitina.mariia@phystech.edu} \\
% 	\And
%  	Роман Исаченко \\
% 	\texttt{roman.isachenko@phystech.edu}
% }

\author{
  Задворнов Егор Валерьевич \\
  МФТИ \\
  \texttt{zadvornov.ev@phystech.edu} \\
\date{\today}
}

\usepackage{setspace}
%\полуторный интервал
\onehalfspacing

\begin{document}

%-----------------

%-----------------



\newgeometry{top=20mm,bottom=20mm,left=30mm,right=15mm}
% \maketitle
% \input{title}
\setcounter{page}{2}

\begin{abstract}
Исследуется проблема прогнозирования распределения тематик в медиапотоке. Показано, что классические регрессионные подходы прогнозирования неустойчивы к появлению нерегулярных циклических циклических изменений с переменным периодом. Предлагается метод ARFilter, который повышает точноть регрессионных моделей с помощью выделения подмножества слов словаря, коррелирующих с временным распределением соответствующих тематик. Для отобранных слов строятся локальные прогнозы их нормированных долей, которые затем агрегируются в прогноз распределения тематик. На синтетических данных метод снижает ошибку прогноза в десять раз по сравнению с базовым подходом ARIMA. Разработанный подход не зависит от выбора конкретной авторегрессионной модели и эффективен для прогнозирования высоковолатильных тематических распределений.
\end{abstract}

\keywords{Временные ряды \and Волатильность \and Корреляционный фильтр \and Медиапоток \and Тематическое моделирование \and ARIMA \and Разложение Дуба.}



\newpage
\tableofcontents
\clearpage % start the table of contents on a new page


% ------------------------------------------
\newpage
\section{Введение}
Способность точно прогнозировать тематические распределения в медиапотоках представляет собой критически важную задачу для современных приложений в маркетинге, медиапроизводстве и исследованиях инноваций. Несмотря на значительный прогресс в методах анализа временных рядов и тематического моделирования, существующие подходы демонстрируют ограниченную эффективность при работе с высоковолатильными текстовыми данными, где присутствие шумовых слов существенно снижает точность прогнозов. Данная работа предлагает метод ARFilter, сочетающий тематическое моделирование с динамической фильтрацией волатильных компонент, что обеспечивает снижение среднеквадратичной ошибки на 91\% относительно классических подходов при сохранении интерпретируемости результатов.

Ключевыми для понимания работы являются следующие понятия: \textit{медиапоток} как хронологическая последовательность документов \(\mathcal{D} = \{(d_i,t_i)\}_{i=1}^{N}\), \textit{словарь} \(\mathcal{W}_{\mathrm{all}}\) — множество \(n\)-грамм (\(1 \leq n \leq 4\)), \textit{тематика} \(T_k \subset \mathcal{W}_{\mathrm{all}}\) — семантический кластер слов, образующий разбиение словаря, и \textit{доля темы} \(\nu_k(t) = \frac{f(T_k,t)}{\sum_j f(T_j,t)}\), где \(f(w,t)\) — частота слова в момент времени \(t\). Тематическое моделирование понимается как метод выявления скрытой тематической структуры в текстовых корпусах.

Основная идея метода заключается в комбинации динамического отбора слов по их корреляции с будущим распределением тематик и иерархического прогнозирования нормированных частот отобранных core-слов с последующей агрегацией на уровне тематик. Этот подход основан на фундаментальных принципах причинного анализа Грангера \cite{Granger1969}, обеспечивающего теоретическую базу для отбора информативных признаков, иерархического прогнозирования \cite{Hyndman2011}, демонстрирующего эффективность композитных моделей, и теории разложения стохастических процессов Дуба \cite{Doob1953}, предоставляющей математический аппарат для разделения детерминированных и стохастических компонент временных рядов \cite{Hamilton1994} \cite{Engle1982}. 

Связь с разложением Дуба: 
В контексте нашей работы \textbf{ядро метода ARFilter} (\(\mathcal{W}_{\text{core}}\)) соответствует \textit{предсказуемой компоненте} в разложении Дуба, характеризующей устойчивые тренды и сезонные паттерны частот слов. \textbf{Шумовые компоненты} (\(\mathcal{W}_{\text{WN}} \cup \mathcal{W}_{\text{SD}}\)) аналогичны \textit{мартингальной составляющей}, описывающей непредсказуемые флуктуации с нулевым математическим ожиданием приращений. Теорема Дуба формализует принцип декомпозиции процесса \(X_t = A_t + M_t\), где:
- \(A_t\) — предсказуемый процесс (аналог \(\sum_{w \in \mathcal{W}_{\text{core}}} f(w,t)\))
- \(M_t\) — мартингал (аналог \(\sum_{w \in \mathcal{W}_{\text{noise}}} f(w,t)\)) \cite{Engle1982} 


Развитие этих идей прослеживается в работах по тематическому моделированию \cite{blei2003latent, sinaga2020unsupervised}, 
обеспечивших переход от анализа отдельных слов к семантическим кластерам, и исследованиях волатильности медиаданных 
\cite{rappaz2019dynamic}, выявивших проблему шумовых слов. 
\textbf{Фундаментальный вклад в развитие методов внесли работы по адаптивной фильтрации} \cite{Cichocki2009}, 
\textbf{предложившие новые подходы к обработке сигналов в шумной среде, и исследования динамических тематических моделей} 
\cite{Blei2012}, \textbf{расширившие возможности анализа эволюции тем во времени}. 
Параллельно развитие методов анализа временных рядов \cite{Shumway2000, Bouchaud2018, Tsay2005, Shumway2017} 
показало потенциал комплексных подходов, однако их применение к тематическим распределениям оставалось неисследованным, 
а перенос принципов фильтрации \cite{Gardner1985} на высокоразмерные текстовые данные требовал адаптации 
к специфике композиционных рядов. \textbf{Современные руководства по прогнозированию} \cite{Hyndman2018} 
\textbf{подчёркивают важность иерархических методов, но не решают проблему шумовых компонент в тематических распределениях}.

Современное состояние области характеризуется рядом ограничений. Традиционные методы прогнозирования 
(ARIMA \cite{box1970time}, экспоненциальное сглаживание \cite{Gardner1985}, Prophet \cite{Taylor2018}) 
работают с единичными рядами и чувствительны к волатильности медиаданных \cite{rappaz2019dynamic}. 
Подходы на основе нейронных сетей \cite{Das2020VAE} и вариационных автокодировщиков теряют интерпретируемость 
и требуют больших объемов данных. Методы кластеризации временных рядов \cite{Grabovoy2019} и тематического моделирования 
\cite{sinaga2020unsupervised, blei2003latent} решают частичные задачи, но не обеспечивают комплексного прогноза распределений. 
\textbf{Моделирование финансовых временных рядов} \cite{Zivot2003} \textbf{предлагает продвинутые техники для волатильных данных, 
но не адаптировано к семантическим особенностям медиапотоков}. 
Попытки прогнозирования трендов \cite{choi2012predicting, holt2004forecasting} либо игнорируют тематический контекст, 
либо используют статическую фильтрацию, неадаптивную к изменению структуры шума, что особенно проблематично 
при наличии white-noise и spike-decay слов, вызывающих рост дисперсии и падение автокорреляции тематических рядов 
(\(\rho_k(1) < 0.3\)).

Для преодоления этих ограничений предлагается метод ARFilter, вводящий два ключевых механизма: \textit{корреляционный фильтр} \(\theta_w = \mathbb{I}[\operatorname{Corr}(f(w,t), \nu_k(t+1)) \geq r]\) для динамического отбора слов и \textit{иерархическую агрегацию} \(\widehat{\nu}_k(t+1) = \sum_{w \in T_k} \theta_w g_{\psi_w}(s(w,t))\) для построения прогноза. Преимущества метода включают \textbf{точность} (снижение MSE с \(23.0 \times 10^{-3}\) до \(2.0 \times 10^{-3}\) относительно ARIMA) и \textbf{совместимость} с различными AR-моделями и кластеризаторами (LDA, K-Means). Основное ограничение — необходимость кросс-валидации для выбора порога \(r\) — компенсируется теоретическими гарантиями (Теорема Задворнова) и экспериментальной эффективностью.

Научная новизна работы заключается в трех аспектах: \textit{новый критерий отбора} слов по корреляции с будущим распределением тематик (а не с собственным будущим), \textit{теоретическое доказательство} снижения MSE при фильтрации волатильных слов, и \textit{адаптация} аппарата разложения Дуба к задаче фильтрации семантических шумов в композиционных временных рядах. Потенциальное влияние включает приложения в маркетинге для сегментированных аудиторий, социологии для анализа общественных трендов и наукометрии для прогнозирования научных направлений, обеспечивая универсальный инструмент для работы с высоковолатильными распределениями в динамически меняющихся средах.





\newpage
\section{Постановка задачи прогнозирования тематических рядов}

Рассмотрим конечное множество тематик $\overline{T}=\{T_1,\dots,T_K\}$ и словарь $\mathcal{W}_{\mathrm{all}}$ из $W=|\mathcal{W}_{\mathrm{all}}|$ слов; каждому кластеру $T_k$ сопоставляется подмножество слов $T_k\subset\mathcal{W}_{\mathrm{all}}$ так, что $T_i\cap T_j=\varnothing$ при $i\neq j$ и $\bigcup_{k=1}^K T_k=\mathcal{W}_{\mathrm{all}}$. Для каждого слова $w\in\mathcal{W}_{\mathrm{all}}$ и момента времени $t\in\{0,\dots,T-1\}$ задана частота $f(w,t)\in\mathbb{N}_{\ge0}$, определяющая число вхождений слова $w$ в тематическом потоке текстов; агрегированная частота темы $T_k$ равна
\[
f(T_k,t)=\sum_{w\in T_k}f(w,t),
\]
а распределение тематик во времени задаётся
\[
\nu_k(t)=\frac{\sum_{w\in T_k}f(w,t)}{\sum_{j=1}^K\sum_{w\in T_j}f(w,t)},\qquad\overline{\nu}(t)=(\nu_1(t),\dots,\nu_K(t)),\qquad\sum_{k=1}^K\nu_k(t)=1.
\]
Предполагается существование оценки $\widehat{\overline{\nu}}\colon\mathbb{N}^{W\times T}\to\mathbb{R}^K$, удовлетворяющей
\[
\widehat{\overline{\nu}}(t)=\widehat{\overline{\nu}}\bigl(t\mid\overline{\phi},\,t-1,\dots,t-t_0\bigr),\qquad\sum_{k=1}^K\widehat{\nu}_k(t)=1,\qquad\mathbb{E}\bigl[\|\overline{\nu}(t)-\widehat{\overline{\nu}}(t)\|_2^2\bigr]\to\min_{\overline{\phi}}.
\]
Данные обладают следующими свойствами: (i) $f(w,t)\in\mathbb{N}_{\ge0}$ для всех $w$ и $t$; (ii) формула для $\nu_k(t)$ задаёт корректную нормировку $\sum_{k}\nu_k(t)=1$; (iii) при малых $f(w,t)\ll1$ дисперсия $\operatorname{Var}[f(w,t)]$ может быть существенной. В качестве модели выступает прогноз распределения тематик на момент $t$ по предыстории длиной $t_0$
\[
\widehat{\overline{\nu}}\!\bigl(t\,\bigm|\,\overline{\phi},\,t{-}1{:}t{-}t_0\bigr),\qquad\overline{\phi}\in\mathbb{R}^d,
\]
а его качество измеряется среднеквадратичной ошибкой
\[
\mathcal{L}(\overline{\phi})=\sum_{t}\bigl\|\overline{\nu}(t)-\widehat{\overline{\nu}}(t\mid\overline{\phi})\bigr\|_2^2
\]
с последующей минимизацией
\[
\min_{\overline{\phi}\in\mathbb{R}^d}\sum_{t}\Bigl\|\overline{\nu}(t)-\widehat{\overline{\nu}}\!\bigl(t\,\bigm|\,\overline{\phi},\,t{-}1{:}t{-}t_0\bigr)\Bigr\|_2^2.
\] \label{eq:classic_min}
Для проверки обобщающей способности используется стратегия скользящего окна: интервал $\{0,\dots,T-1\}$ разбивается на обучающий и тестовый сегменты длинами $T_{\mathrm{train}}$ и $T_{\mathrm{test}}$, причём на каждом шаге модель обучается на отрезке $[t{-}t_0,\,t{-}1]$ и проверяется в точке $t$. Итоговая цель — построить устойчивую модель, способную надёжно прогнозировать распределение тематических частот в присутствии аномалий и цикличных колебаний.


% ------------------------------------------

\newpage
\section{Построение тематических кластеров}\label{ssec:cluster_build}
Пусть $\mathcal{W}_{\mathrm{all}}=\{w_1,\dots,w_W\}$ — словарь $n$-грамм ($1\le n\le4$). Требуется построить разбиение $\overline{T}=\{T_1,\dots,T_K\}$ такое, что $T_i\cap T_j=\varnothing$ и $\bigcup_{k=1}^{K}T_k=\mathcal{W}_{\mathrm{all}}$. 

\paragraph{Метод~1 (Embeddings+$k$-means).}
Определяется отображение $\mathrm{E}\colon\mathcal{W}_{\mathrm{all}}\to\mathbb{R}^d$, $\mathrm{E}(w)=\mathbf e_w$, получающее контекстуальные встраивания слов при помощи модели \texttt{all-MiniLM-L12-v2} из библиотеки Sentence Transformers \cite{Reimers2019,SentenceTransformers}. Для фиксированного числа кластеров $K'$ решается задача 
\[
\min_{\{T_k\}_{k=1}^{K'}}\sum_{k=1}^{K'}\sum_{w\in T_k}\bigl\|\mathbf e_w-\boldsymbol\mu_k\bigr\|_2^{2},\qquad
\boldsymbol\mu_k=\frac{1}{|T_k|}\sum_{w\in T_k}\mathbf e_w,
\]
то есть применяется алгоритм $k$-средних к полученной матрице встраиваний; подробности алгоритма приводятся в \cite{sinaga2020unsupervised}, поэтому здесь они опущены. В результате каждое слово $w\in\mathcal{W}_{\mathrm{all}}$ однозначно попадает в один из кластеров $T_k$, удовлетворяя исходным условиям разбиения.

\paragraph{Метод~2 (Latent Dirichlet Allocation).}
Второй алгоритм работает следующим образом. Этапы предварительной обработки включают лемматизацию, удаление стоп-слов и отбор слов длиной от двух до трёх символов, после чего обучается модель латентного размещения Дирихле (Latent Dirichlet Allocation, LDA) \cite{blei2003latent}. LDA представляет собой вероятностную модель тем, выявляющую скрытые (латентные) темы в корпусе документов; она исходит из предположения, что каждый документ является смесью нескольких тем, а каждое слово принадлежит одной из этих тем.

Алгоритм LDA формулируется так: сначала для каждой темы $k\in\{1,\ldots,K'\}$ выбирается распределение слов $\phi_k\sim\mathrm{Dirichlet}(\beta)$, где $\beta$ служит гиперпараметром априорного распределения Дирихле; далее для каждого документа $d\in\{1,\ldots,D\}$ выбирается распределение тем $\theta_d\sim\mathrm{Dirichlet}(\alpha)$ с гиперпараметром $\alpha$, после чего для каждого слова $w_{dn}$ внутри документа $d$ последовательно выбираются тема $z_{dn}\sim\mathrm{Multinomial}(\theta_d)$ и само слово $w_{dn}\sim\mathrm{Multinomial}(\phi_{z_{dn}})$. Совместное распределение всех переменных модели — распределений тем $\theta$, распределений слов $\phi$, назначенных тем $\mathbf z$ и наблюдаемых слов $\mathbf w$ — записывается следующим образом:
\begin{equation}
p(\theta, \phi, \mathbf{z}, \mathbf{w} \mid \alpha, \beta) = \prod_{k=1}^{K'} p(\phi_k \mid \beta) \prod_{d=1}^D p(\theta_d \mid \alpha) \prod_{n=1}^{N_d} p(z_{dn} \mid \theta_d)\, p(w_{dn} \mid \phi_{z_{dn}})
\end{equation}
Задача вывода в модели LDA заключается в вычислении апостериорного распределения скрытых переменных $\theta$, $\phi$ и $\mathbf z$ по наблюдаемым словам $\mathbf w$ при заданных гиперпараметрах $\alpha$ и $\beta$; из-за вычислительной сложности точный вывод заменяется приближёнными техниками (например, вариационным выводом или коллапсированным сэмплированием Гиббса). Кластеризация затем осуществляется правилом \(w\mapsto T_{\arg\max_{k}\theta_{w,k}}.\)

Оптимальное значение $K$ определяется по метрике средней когерентности $C_{\text{cv}}$ \cite{roder2015exploring}, измеряющей семантическое сходство слов внутри кластера и, следовательно, его интерпретируемость; выбирается такое $K$, которое минимизирует
\[
K=\arg\min_{K'\in\mathbb N}\frac{1}{K'}\sum_{m=0}^{K'-1}C_{\text{cv}}(c_m).
\]
Мера когерентности $C_{\text{cv}}$ вычисляется как
\[
C_{\text{cv}}=\frac{1}{\lvert S^{\text{one}}_{\text{set}}\rvert}\sum_{(W_0,W_*)\in S^{\text{one}}_{\text{set}}}\tilde m_{\cos(\mathrm{nlr},1)}(W_0,W_*),
\]
где $S^{\text{one}}_{\text{set}}=\{(W_0,W_*)\mid W_0=\{w_i\},\,w_i\in W,\,W_*=W\}$ — сегментация, задающая множество пар подмножеств слов, а $\tilde m_{\cos(\mathrm{nlr},1)}(W_0,W_*)$ — косинусное сходство между контекстными векторами $W_0$ и $W_*$, измеряющее, насколько хорошо $W_*$ «объясняет» $W_0$.

\paragraph{Выбор метода.}
Цель данного этапа состоит в сравнении описанных методов по качеству кластеризации и выборе одного из них для дальнейшего использования в задаче прогнозирования; для этого введён критерий, основанный на когерентности: пусть $(C_{cv}^{\mathrm{Emb}},C_{cv}^{\mathrm{LDA}})$ обозначает значения когерентности при оптимальном числе кластеров $K^{\star}$, тогда
\[
\text{Модель}= 
\begin{cases}
\text{Embeddings+$k$-means}, & |C_{cv}^{\mathrm{Emb}}-C_{cv}^{\mathrm{LDA}}|<0.1\ \text{и}\ C_{cv}^{\mathrm{Emb}}>0.5,\\[4pt]
\text{LDA}, & \text{иначе},
\end{cases}
\]
что допускает небольшое снижение метрики у $k$-means по сравнению с LDA, поскольку первый подход проще и выигрывает по ряду иных характеристик; далее, для каждого момента времени $t\in\{0,\dots,T-1\}$ и каждой темы $T_k\in\overline{T}$ определяются агрегированные частоты $f(T_k,t)=\sum_{w\in T_k}f(w,t)$ и относительные значения
\[
\nu_k(t)=\frac{f(T_k,t)}{\sum_{j=1}^{K}f(T_j,t)},\label{eq:nu_def}
\]
в результате чего формируются ряды $\{f(T_k,t)\}_t$, служащие исходными данными для последующих прогнозных моделей.

% ----------------------------------
\newpage
\section{Модель прогноза будущих частот}
\subsection{Классический авторегрессионный подход}\label{ssec:classic}
Классический метод основывается на авторегрессионной модели, прогнозирующей частоты тематик по их прошлым значениям; обычно применяются модели семейства ARIMA, например реализация Prophet. Для каждой темы $T_k\in\overline{T}$ вводится скалярный ряд $y_k(t)=\nu_k(t)$, $t\in\{0,\dots,T-1\}$, где $\nu_k(t)$ определена в~(\ref{eq:nu_def}).

% \paragraph{SARIMA$(p_k,d_k,q_k)$ для долей.}
Пусть $B$ — оператор сдвига, $B^j y_k(t)=y_k(t-j)$. Сезонная авторегрессионная интегрированная скользящая средняя $SARIMA(p,d,q)(P,D,Q)_m$, являющаяся расширением ARIMA с учётом сезонности, описывается уравнением
\[
(1-\phi_1 B)(1-\Phi_1 B^m)(1-B)^d(1-B^m)^D y_t = (1+\theta_1 B)(1+\Theta_1 B^m)\epsilon_t,
\]
где $p$ - порядок несезонной авторегрессии, $d$ - порядок несезонного дифференцирования, $q$ - порядок несезонного скользящего среднего, $P$ - порядок сезонной авторегрессии, $D$ - порядок сезонного дифференцирования, $Q$ - порядок сезонного скользящего среднего, $m$ - количество периодов в сезонном цикле, $\phi_1$ - параметр несезонной авторегрессии, $\Phi_1$ - параметр сезонной авторегрессии, $\theta_1$ - параметр несезонного скользящего среднего, $\Theta_1$ - параметр сезонного скользящего среднего, $B$ - оператор сдвига назад ($B^j y_t = y_{t-j}$), а $\epsilon_t$ - белый шум.

Алгоритм~SARIMA реализуется последовательно: сперва выполняется идентификация, то есть выбор порядков $p$, $d$, $q$, $P$, $D$, $Q$ и длины сезонного цикла $m$ на основе графиков автокорреляционной (ACF) и частичной автокорреляционной (PACF) функций и тестов на стационарность; затем происходит оценка параметров $\phi_1$, $\Phi_1$, $\theta_1$, $\Theta_1$ и дисперсии шума $\sigma^{2}$, обычно методом максимального правдоподобия или условных наименьших квадратов; после этого проводится диагностика адекватности модели через анализ остатков на автокорреляцию, нормальность распределения и гомоскедастичность; и, наконец, формируются точечные и интервальные прогнозы по рекуррентной схеме
\[
\hat{y}_{t+h\mid t}= \phi_1\hat{y}_{t+h-1\mid t}+ \Phi_1\hat{y}_{t+h-m\mid t}-\theta_1\hat{\epsilon}_{t+h-1\mid t}-\Theta_1\hat{\epsilon}_{t+h-m\mid t},
\qquad \hat{\epsilon}_{t+h\mid t}=0\ \text{при}\ h>0,
\]
где $\hat{y}_{t+h\mid t}$ — прогноз на горизонт $h$ шагов вперёд, полученный на основе оценённой модели.

Точечный прогноз определяется как
\[
\widehat{y}_k(t+1)=\mathbb{E}_{\widehat{\psi}_k}\bigl[y_k(t+1)\mid\mathcal F_t\bigr],
\]
а для обеспечения нормировки $\sum_{k=1}^{K}\widehat{\nu}_k(t+1)=1$ применяется преобразование
\[
\widehat{\nu}_k(t+1)=\frac{\widehat{y}_k(t+1)}{\sum_{j=1}^{K}\widehat{y}_j(t+1)},
\]
при этом качество прогноза оценивается метриками MAE и MSE, что эквивалентно функционалу
\[
\mathcal{L}=\sum_{t}\|\overline{\nu}(t)-\widehat{\overline{\nu}}(t)\|_2^2.
\]

Подбор $(p_k,d_k,q_k)$ и вычисление $\widehat{y}_k(t+1)$ выполнялись с помощью библиотеки \texttt{Prophet} \cite{prophet}, представляющей временной ряд в виде аддитивной регрессии
\[
y(t)=g(t)+s(t)+h(t)+\epsilon_t,
\]
где $g(t)$ отвечает за тренд, $s(t)$ за сезонность, $h(t)$ за праздничные эффекты, а $\epsilon_t$ моделирует случайную ошибку. Тренд $g(t)$ может задаваться кусочно-линейной зависимостью
\[
g_{\text{lin}}(t)=(k+\sum_{j=1}^{S}\delta_j)t+(m+\sum_{j=1}^{S}\gamma_j)
\]
или логистическим ростом
\[
g_{\text{log}}(t)=\frac{C(t)}{1+\exp\!\bigl(-g_{\text{lin}}(t)\bigr)},
\]
что позволяет гибко описывать как линейные, так и насыщаемые динамики.

Сезонная компонента описывается рядом Фурье 
\[
s(t)=\sum_{n=1}^{N}\bigl(a_n\cos(2\pi n t/P)+b_n\sin(2\pi n t/P)\bigr),
\]
праздничные эффекты моделируются суммой индикаторных функций
\[
h(t)=\sum_{i=1}^{L}\theta_i\mathbb 1_{t\in D_i},
\]
а параметры модели оцениваются максимизацией правдоподобия, что эквивалентно решению
\(
\hat{\theta}=\arg\min_{\theta}\Bigl[-\log p(y\mid\theta)\Bigr]=\arg\min_{\theta}\left\{\tfrac{n}{2}\log(2\pi)+\tfrac{n}{2}\log(\sigma^{2})+\tfrac{1}{2\sigma^{2}}\sum_{t}\bigl(y(t)-\sum_{i}f_i(t)\bigr)^{2}\right\},
\)
где $\theta$ — вектор параметров, $n$ — число наблюдений, $\sigma^{2}$ — дисперсия ошибки, $y(t)$ — фактические значения ряда, а $f_i(t)$ — отдельные компоненты модели.

Прогноз формируется экстраполяцией оценённого тренда, добавлением сезонной составляющей и учётом праздничных эффектов, что даёт
\[
\hat{y}(t+h)=\hat{g}(t+h)+\hat{s}(t+h)+\hat{h}(t+h).
\]

% ----------------------------------------------

\newpage
\subsection{Проблемы и ограничения классического подхода}\label{ssec:limits}
Основная проблема классического подхода заключается в том, что он усредняет все слова с равными весами, не учитывая их различную предсказательную способность относительно будущего распределения тематик. Анализ временных рядов частот слов показывает, что слова в словаре могут иметь принципиально разную временную структуру и, следовательно, разный вклад в точность прогноза. Поэтому на основании корреляционных свойств и характера временной эволюции словарь целесообразно разложить на три непересекающиеся части 
\[
\mathcal{W}_{\mathrm{all}}=\mathcal{W}_{\mathrm{core}}\cup\mathcal{W}_{\mathrm{WN}}\cup\mathcal{W}_{\mathrm{SD}},\qquad \mathcal{W}_{a}\cap\mathcal{W}_{b}=\varnothing,\;a\neq b,
\]
при этом ядро $\mathcal{W}_{\mathrm{core}}$ составляют слова с устойчивой сезонно-трендовой структурой и высокой корреляцией с будущими значениями соответствующих тематик,
\[
\mathcal{W}_{\mathrm{core}}=\bigl\{w\in\mathcal{W}_{\mathrm{all}}\mid \rho_w:=\operatorname{Corr}\bigl(f(w,t),f(T_k,t+1)\bigr)\ge r,\;w\in T_k\bigr\},
\]
а именно эти слова обладают наибольшей предсказательной ценностью и формируют основу устойчивой динамики тематических рядов.

Вторую группу образуют white-noise слова $\mathcal{W}_{\mathrm{WN}}$, то есть мультитематические слова, которые часто переходят между темами и потому моделируются случайной ошибкой без предсказательной силы; их частоты описываются
\[
\mathcal{W}_{\mathrm{WN}} = \left\{ w \in \mathcal{W}_{\mathrm{all}} \;\middle|\; f(w,t) = [\varepsilon_w(t)]_+,\, \rho_w = 0 \right\}
\]
где $\varepsilon_w(t)\overset{\text{i.i.d.}}{\sim}\mathcal N(\mu_{\mathrm{WN}},\sigma_{\mathrm{WN}}^{2})$, $[x]_+=\max\{0,x\}$.
Так что их вклад в $\nu_k(t)$ представляет собой чистый шум и лишь повышает дисперсию прогноза. 

Третью группу составляют spike-decay слова $\mathcal{W}_{\mathrm{SD}}$ - слова с нерегулярными всплесками активности и переменным периодом. Хотя они и сохраняют стационарность ряда, но значительно увеличивают дисперсию ошибки прогноза: 
\[
\mathcal{W}_{\mathrm{SD}} = \left\{ w \in \mathcal{W}_{\mathrm{all}} \;\middle|\; f(w,t) = [S_w(t) + \varepsilon_w(t)]_+, \quad \rho_w(\tau) \approx 0\right\}
\]
где $S_w(t) = \textstyle\sum A_{w,m} \kappa(t{-}\tau_{w,m})$, $\kappa(\tau)=\frac{\mathrm e^{-\lambda|\tau|}}{1+\tau^{2}}$ и $\varepsilon_w(t)\sim\mathcal N(\mu_{\mathrm{SD}},\sigma_{\mathrm{SD}}^{2})$. Сигнал $S_w(t)$ порождает редкие пики с экспоненциальным затуханием; поскольку $\operatorname{Var}f(w,t)\gg \operatorname{Var}f(u,t)$ для $u\in\mathcal{W}_{\mathrm{core}}$, включение таких слов без учёта их особенностей существенно ухудшает качество прогноза.


\paragraph{Воздействие на агрегаты.}  
Для темы \(T_k\) представим
\(
  f(T_k,t)
  \;=\;
  f^{(\mathrm{core})}(T_k,t)
  \;+\;
  f^{(\mathrm{WN})}(T_k,t)
  \;+\;
  f^{(\mathrm{SD})}(T_k,t),
\)
где компоненты соответствуют разбиению слов \(\mathcal{W}_{\mathrm{all}} = \mathcal{W}_{\mathrm{core}} \cup \mathcal{W}_{\mathrm{WN}} \cup \mathcal{W}_{\mathrm{SD}}\). 
Введём обозначения:
\(
  \sigma_{\mathrm{WN}}^{2} \;:=\; \operatorname{Var}\bigl(f^{(\mathrm{WN})}(T_k,t)\bigr),
  \qquad
  \sigma_{\mathrm{SD}}^{2} \;:=\; \operatorname{Var}\bigl(f^{(\mathrm{SD})}(T_k,t)\bigr).
\)

\begin{theorem}{Утверждение 1 (Задворнов, 2025)}
Пусть \(
\mathcal{W}_{\mathrm{all}} = \mathcal{W}_{\mathrm{core}} \cup \mathcal{W}_{\mathrm{WN}} \cup \mathcal{W}_{\mathrm{SD}}, \quad \mathcal{W}_a \cap \mathcal{W}_b = \varnothing
\)
при условии попарной некоррелированности компонент
\[
\forall\,a \neq b,\;\forall\,\tau \in \mathbb{Z}:\quad
Cov\bigl(f^{(a)}(T_k,t),\,f^{(b)}(T_k,t+\tau)\bigr) = 0
\]
Тогда
\[
Cov\bigl(f(T_k,t),\,f(T_k,t+1)\bigr) = Cov\bigl(f^{(\mathrm{core})}(t),\,f^{(\mathrm{core})}(t{+}1)\bigr)
\]
и
\[
\rho_k(1) = \frac{Cov^{\mathrm{core}}}{Var^{\mathrm{core}} + \sigma_{\mathrm{WN}}^2 + \sigma_{\mathrm{SD}}^2}.
\]
\end{theorem}

\begin{proof}
\textbf{Доказательство п.1:}
Распишем ковариацию, используя линейность и условие некоррелированности:
\[
\begin{aligned}
&\operatorname{Cov}\bigl(f(T_k,t), f(T_k,t+1)\bigr) \\
&= \sum_{a,b} \operatorname{Cov}\bigl(f^{(a)}(t), f^{(b)}(t+1)\bigr) \\
&= \sum_{a} \operatorname{Cov}\bigl(f^{(a)}(t), f^{(a)}(t+1)\bigr) \quad \text{(при \(a \neq b\) ковариации нулевые)}.
\end{aligned}
\]
Для \(\mathcal{W}_{\mathrm{WN}}\) по определению \(\rho_w = 0\), откуда \(\operatorname{Cov}\bigl(f^{(\mathrm{WN})}(t), f^{(\mathrm{WN})}(t+1)\bigr) = 0\). \\
Для \(\mathcal{W}_{\mathrm{SD}}\) по условию \(\rho_w(\tau) \approx 0\) при \(\tau \neq 0\), следовательно, \(\operatorname{Cov}\bigl(f^{(\mathrm{SD})}(t), f^{(\mathrm{SD})}(t+1)\bigr) \approx 0\). \\
Таким образом:
\[
\operatorname{Cov}\bigl(f(T_k,t), f(T_k,t+1)\bigr) = \operatorname{Cov}\bigl(f^{(\mathrm{core})}(t), f^{(\mathrm{core})}(t+1)\bigr).
\]

\textbf{Доказательство п.2:}
Дисперсия суммы:
\[
\begin{aligned}
\operatorname{Var} f(T_k,t) 
&= \operatorname{Var}\left( \sum_{a} f^{(a)}(t) \right) \\
&= \sum_{a} \operatorname{Var}\bigl(f^{(a)}(t)\bigr) \quad \text{(ковариации между компонентами нулевые)} \\
&= \operatorname{Var} f^{(\mathrm{core})}(t) + \sigma_{\mathrm{WN}}^{2} + \sigma_{\mathrm{SD}}^{2}.
\end{aligned}
\]
Лаг-1 автокорреляция:
\[
\rho_k(1) = \frac{\operatorname{Cov}\bigl(f(t), f(t+1)\bigr)}{\operatorname{Var} f(t)} = \frac{\operatorname{Cov}^{(\mathrm{core})}}{\operatorname{Var} f^{(\mathrm{core})}(t) + \sigma_{\mathrm{WN}}^{2} + \sigma_{\mathrm{SD}}^{2}}.
\]
\end{proof}

Из утверждения 1 непосредственно следует:
\[
\rho_{k}(1)
\;=\;
\frac{\operatorname{Cov}\bigl(f^{(\mathrm{core})}(T_k,t),\,f^{(\mathrm{core})}(T_k,t+1)\bigr)}
     {\operatorname{Var} f(T_k,t)}
\;\le\;
\rho_{k}^{\mathrm{core}}
\,\frac{\operatorname{Var} f^{(\mathrm{core})}(T_k,t)}
       {\operatorname{Var} f^{(\mathrm{core})}(T_k,t)+\sigma_{\mathrm{WN}}^{2}+\sigma_{\mathrm{SD}}^{2}},
\]
где \(\rho_{k}^{\mathrm{core}}\) – лаг-1 автокорреляция ядра.

Вывод для ARIMA-прогноза: чем больше $\sigma_{\mathrm{WN}}^{2}$ и $\sigma_{\mathrm{SD}}^{2}$, тем сильнее
снижается $\rho_{k}(1)$, нарушая предпосылку AR-модели о существенной
корреляционной зависимости, и тем сильнее размывается сигнал, что приводит к деградации AR-прогноза.
Отсюда необходимость альтернативного подхода, который учётывает $\mathcal{W}_{\mathrm{WN}}\cup\mathcal{W}_{\mathrm{SD}}$ в прогностической модели.


% ----------------------------------------------

\newpage
\subsection{Альтернативный корреляционный подход}
\label{ssec:alt}

Альтернативный подход заключается в рассмотрении вектора корреляций частот слов с будущими значениями их тематик (корреляционный метод) и выделении для прогноза тематики только слов с наибольшими корреляционными значениями.
Далее для каждого слова $w_k$ построить отдельную AR-модель, предсказывающую популярность ( относительную долю ) слова $w_k$ среди core-слов $\mathcal{W}_{\mathrm{core}}$ в момент $t$.
Относительная доля считается как количество упоминаний слова за период, поделённое на количество упоминаний выделенных слов за период.
Далее суммировать и нормировать прогнозы слов, чтобы в итоге получить популярность тематики как сумму популярностей слов внутри неё.

В рамках корреляционного скрининга для каждого $w\in\mathcal{W}_{\mathrm{all}}$ оценивается корреляция с будущим значением соответствующей тематики 
\[
\widehat\rho_{w}:=\operatorname{Corr}\bigl(f(w,t),f(T_{k},t+1)\bigr),\qquad t\in T_{\mathrm{train}},\;w\in T_{k},
\]
после чего вводится индикатор 
\[
\theta_{w}=\mathbb{I}\bigl[\widehat\rho_{w}\ge r\bigr],\qquad r\in(0,1),\qquad\overline{\theta}\in\{0,1\}^{W},\qquad 
\mathcal{W}_{\mathrm{core}}=\{w:\theta_{w}=1\},
\]
и для $t\in\{0,\dots,T-1\}$ формируются нормированные ряды
\[
s(w,t)=\frac{f(w,t)}{\displaystyle\sum_{v\in\mathcal{W}_{\mathrm{core}}}f(v,t)},\qquad w\in\mathcal{W}_{\mathrm{core}}.
\]
Для каждого core-слова $w\in\mathcal{W}_{\mathrm{core}}$ формируется авторегрессионный прогноз $\widehat{s}\!\bigl(w,t+1\mid\psi_{w}\bigr)=g_{\psi_{w}}\bigl(s(w,t),\dots,s(w,t-t_{0}+1)\bigr)$, где вектор параметров $\psi_{w}\in\mathbb{R}^{d_{w}}$ оценивается на обучающем интервале методом наименьших квадратов; далее все полученные значения переводятся в вероятностную шкалу по правилу
\[
\widetilde{s}(w,t+1)=\frac{\widehat{s}(w,t+1\mid\psi_{w})}{\sum_{v\in\mathcal{W}_{\mathrm{core}}}\widehat{s}(v,t+1\mid\psi_{v})},\qquad \sum_{w\in\mathcal{W}_{\mathrm{core}}}\widetilde{s}(w,t+1)=1,
\]
после чего относительная частота темы $T_{k}$ в момент $t+1$ вычисляется агрегированием нормированных прогнозов
\[
\widehat{\nu}_{k}(t+1)=\sum_{w\in T_{k}}\theta_{w}\,\widetilde{s}(w,t+1),\qquad\sum_{k=1}^{K}\widehat{\nu}_{k}(t+1)=1,
\]
Таким образом $\mathcal{W}_{\mathrm{WN}}\cup\mathcal{W}_{\mathrm{SD}}$  
автоматически исключаются при $\widehat\rho_{w}<r$,  
а прогноз строится только на информативном подмножестве
$\mathcal{W}_{\mathrm{core}}$.

% \paragraph{Оптимизационная постановка.}
% Разделив параметры на  
% \(
%   \overline{\phi}=(\overline{\theta},\overline{\psi}),
% \;
%   \overline{\psi}=(\psi_{w})_{w:\theta_{w}=1},
% \)
% получаем задачу
% \[
%   \min_{\overline{\theta},\overline{\psi}}
%   \frac{1}{|T_{\mathrm{test}}|}
%   \sum_{t\in T_{\mathrm{test}}}
%   \bigl\|
%     \overline{\nu}(t)-\widehat{\overline{\nu}}(t\mid
%     \overline{\theta},\overline{\psi},t-1{:}t-t_0)
%   \bigr\|_2^{2},
%   \quad
%   \text{s.t. } \theta_{w}=0\;\;\forall w:\widehat\rho_{w}<r.
% \]

\begin{theorem}{Теорема о снижении MSE при фильтрации слов (Задворнов, 2025)}
Пусть для тематики $T_k$:
\begin{enumerate}
    \item $\mu_k = \mathbb{E}[\nu_k(t) \mid \mathcal{F}_{t-1}]$ - условное матожидание распределения по всем словам
    \item $\mu_k^{(\mathrm{core})} = \mathbb{E}[\widehat{\nu}_k^{(\mathrm{core})}(t) \mid \mathcal{F}_{t-1}]$ - условное матожидание по core-словам
    \item Выполнены условия Утверждения 1 о разбиении словаря
\end{enumerate}
и дополнительно выполнены условия на шумовые слова:
\begin{enumerate}
    \item Для $w \in \mathcal{W}_{\mathrm{WN}}$: $\mathbb{E}[f(w,t) \mid \mathcal{F}_{t-1}] \leq C_{\mathrm{WN}} \sigma_{\mathrm{WN}}$ 
    \item Для $w \in \mathcal{W}_{\mathrm{SD}}$: $\mathbb{E}[f(w,t) \mid \mathcal{F}_{t-1}] \leq C_{\mathrm{SD}} \sigma_{\mathrm{SD}}$ (в отсутствие пиков)
    \item $|\mathcal{W}_{\mathrm{WN}}| C_{\mathrm{WN}}^2 + |\mathcal{W}_{\mathrm{SD}}| C_{\mathrm{SD}}^2 < \kappa$
\end{enumerate}
где $\kappa > 0$ - минимальное собственное значение информационной матрицы модели прогноза.

Тогда:
\[
  \mathbb{E}\!\bigl[\mathcal L_{\mathrm{ARFilter}}\mid \mathcal{F}_{t}\bigr]:=\mathrm{MSE}\bigl(\widehat{\nu}_k^{(\mathrm{core})}(t)\bigr)
  \;<\; 
  \mathrm{MSE}\bigl(\widehat{\nu}_k^{(\mathrm{all})}(t)\bigr) =:\mathbb{E}\!\bigl[\mathcal L_{\mathrm{baseline}}\mid \mathcal{F}_{t}\bigr]
\]
\end{theorem}

\begin{proof}
\textbf{Шаг 1: Ошибки прогнозов}
\[
{\nu}_k(T-1) = \mu_k + \varepsilon_k, \quad \mathbb{E}[\varepsilon_k \mid \mathcal{F}_{T-2}] = 0, \quad \text{Var}(\varepsilon_k \mid \mathcal{F}_{T-2}) = \sigma_k^2.
\]
Из условия несмещенности для полного прогноза:
\(
\mathbb{E}[\widehat{\nu}_k^{(\mathrm{all})}(T-1) \mid \mathcal{F}_{T-2}] = \mu_k
\)

Для прогноза по стабильным словам:
\[
\mathbb{E}[\widehat{\nu}_k^{(\mathrm{core})}(T-1) \mid \mathcal{F}_{T-2}] = \mu_k^{(core)} = \sum_{w \in \mathcal{W}_{\mathrm{core}} \cap T_k} \mathbb{E}[X_w(T-1) \mid \mathcal{F}_{T-2}]
\]
Заметим, что в общем случае \(\mu_k^{(core)} \neq \mu_k\), но нам это не важно для дисперсии.
Ошибки прогнозов:
\[
e^{(\mathrm{all})} = {\nu}_k(T-1) - \widehat{\nu}_k^{(\mathrm{all})}(T-1) = \varepsilon_k + (\mu_k - \widehat{\nu}_k^{(\mathrm{all})}(T-1))
\]
\[
e^{(\mathrm{core})} = \nu_k(T-1) - \widehat{\nu}_k^{(\mathrm{core})}(T-1) = \varepsilon_k + (\mu_k - \widehat{\nu}_k^{(\mathrm{core})}(T-1))
\]
Тогда условное MSE:
\(
\mathbb{E}[(e^{(\mathrm{all})})^2 \mid \mathcal{F}_{T-2}]
= \mathbb{E}[\varepsilon_k^2 \mid \mathcal{F}_{T-2}] + \mathbb{E}[(\mu_k - \widehat{\nu}_k^{(\mathrm{all})}(T-1))^2 \mid \mathcal{F}_{T-2}] + 2 \mathbb{E}[\varepsilon_k (\mu_k - \widehat{\nu}_k^{(\mathrm{all})}(T-1)) \mid \mathcal{F}_{T-2}] = \sigma_k^2 + \text{Var}(\widehat{\nu}_k^{(\mathrm{all})}(T-1) \mid \mathcal{F}_{T-2})
\)
Поскольку $\varepsilon_k$ не зависит от \(\widehat{\nu}_k^{(\mathrm{all})}(T-1)\) (который строится по истории до T-2) и \(\mathbb{E}[\varepsilon_k \mid \mathcal{F}_{T-2}] = 0\), то ковариационное слагаемое равно 0.
Аналогично:
\[
\mathbb{E}[(e^{(\mathrm{core})})^2 \mid \mathcal{F}_{T-2}]
= \sigma_k^2 + \text{Var}(\widehat{\nu}_k^{(\mathrm{core})}(T-1) \mid \mathcal{F}_{T-2}) + (\mu_k - \mu_k^{(core)})^2
\]
Потому что:
\[
e^{(\mathrm{core})} = \varepsilon_k + (\mu_k - \widehat{\nu}_k^{(\mathrm{core})}(T-1))
= \varepsilon_k + (\mu_k - \mu_k^{(core)}) + (\mu_k^{(core)} - \widehat{\nu}_k^{(\mathrm{core})}(T-1))
\]
Тогда:
\[
\mathbb{E}[(e^{(\mathrm{core})})^2 \mid \mathcal{F}_{T-2}]
= \mathbb{E}[\varepsilon_k^2] + (\mu_k - \mu_k^{(core)})^2 + \mathbb{E}[(\mu_k^{(core)} - \widehat{\nu}_k^{(\mathrm{core})}(T-1))^2]
\]
\[
+ 2 \mathbb{E}[\varepsilon_k (\mu_k - \mu_k^{(core)})] + 2 \mathbb{E}[\varepsilon_k (\mu_k^{(core)} - \widehat{\nu}_k^{(\mathrm{core})}(T-1))] + 2 (\mu_k - \mu_k^{(core)}) \mathbb{E}[\mu_k^{(core)} - \widehat{\nu}_k^{(\mathrm{core})}(T-1)]
\]
Заметим:
\(
\mathbb{E}[\varepsilon_k (\mu_k - \mu_k^{(core)}) | \mathcal{F}_{T-2}] = (\mu_k - \mu_k^{(core)}) \mathbb{E}[\varepsilon_k | \mathcal{F}_{T-2}] = 0
\)
\[
\mathbb{E}[\varepsilon_k (\mu_k^{(core)} - \widehat{\nu}_k^{(\mathrm{core})}(T-1)) | \mathcal{F}_{T-2}] = \mathbb{E}[\varepsilon_k | \mathcal{F}_{T-2}] \cdot (\mu_k^{(core)} - \mathbb{E}[\widehat{\nu}_k^{(\mathrm{core})}(T-1) | \mathcal{F}_{T-2}]) = 0
\]

\(\mathbb{E}[\mu_k^{(core)} - \widehat{\nu}_k^{(\mathrm{core})}(T-1) | \mathcal{F}_{T-2}] = 0\)   (по условиям из утрержения 1, примененному к каждому слову и агрегированию)

Поэтому:
\(
\mathbb{E}[(e^{(\mathrm{core})})^2 \mid \mathcal{F}_{T-2}]
= \sigma_k^2 + (\mu_k - \mu_k^{(core)})^2 + \text{Var}(\widehat{\nu}_k^{(\mathrm{core})}(T-1) \mid \mathcal{F}_{T-2})
\)
Итог условное MSE:
\[
\mathrm{MSE}^{(\mathrm{all})} = \mathbb{E} \left[ \left( e^{(\mathrm{all})} \right)^2 \mid \mathcal{F}_{T-2} \right] = \sigma_k^2 + \mathrm{Var}\left( \widehat{\nu}_k^{(\mathrm{all})}(T-1) \mid \mathcal{F}_{T-2} \right)
\]
\[
\mathrm{MSE}^{(\mathrm{core})} = \mathbb{E} \left[ \left( e^{(\mathrm{core})} \right)^2 \mid \mathcal{F}_{T-2} \right] = \sigma_k^2 + \mathrm{Var}\left( \widehat{\nu}_k^{(\mathrm{core})}(T-1) \mid \mathcal{F}_{T-2} \right) + B^2
\]
где $B = \mu_k - \mu_k^{(\mathrm{core})}$, $\sigma_k^2 = \mathrm{Var}(\nu_k(T-1) \mid \mathcal{F}_{T-2})$.

\textbf{Шаг 2: Выражение смещения}
Смещение $B$ равно сумме вкладов шумовых слов:
\[
B = \sum_{w \in \mathcal{W}_{\mathrm{WN}}} \mathbb{E}[f(w,T-1) \mid \mathcal{F}_{T-2}] + \sum_{w \in \mathcal{W}_{\mathrm{SD}}} \mathbb{E}[f(w,T-1) \mid \mathcal{F}_{T-2}]
\]

Используя условия 1 и 2:
\[
|B| \leq \sum_{w \in \mathcal{W}_{\mathrm{WN}}} C_{\mathrm{WN}} \sigma_{\mathrm{WN}} + \sum_{w \in \mathcal{W}_{\mathrm{SD}}} C_{\mathrm{SD}} \sigma_{\mathrm{SD}}
\]
\[
B^2 \leq \left( \sum_{w \in \mathcal{W}_{\mathrm{WN}}} C_{\mathrm{WN}} \sigma_{\mathrm{WN}} + \sum_{w \in \mathcal{W}_{\mathrm{SD}}} C_{\mathrm{SD}} \sigma_{\mathrm{SD}} \right)^2 \leq 2 \left( |\mathcal{W}_{\mathrm{WN}}|^2 C_{\mathrm{WN}}^2 \sigma_{\mathrm{WN}}^2 + |\mathcal{W}_{\mathrm{SD}}|^2 C_{\mathrm{SD}}^2 \sigma_{\mathrm{SD}}^2 \right)
\]

\textbf{Шаг 3: Разность дисперсий прогнозов}
Из Утверждения 1 и свойств моделей ARIMA/Prophet:
\(
\Delta V := \mathrm{Var}\left( \widehat{\nu}_k^{(\mathrm{all})}(T-1) \mid \mathcal{F}_{T-2} \right) - \mathrm{Var}\left( \widehat{\nu}_k^{(\mathrm{core})}(T-1) \mid \mathcal{F}_{T-2} \right) \geq \kappa \left( \sigma_{\mathrm{WN}}^2 |\mathcal{W}_{\mathrm{WN}}| + \sigma_{\mathrm{SD}}^2 |\mathcal{W}_{\mathrm{SD}}| \right)
\)

\textbf{Шаг 4: Сравнение $B^2$ и $\Delta V$}
Покажем, что $B^2 < \Delta V$:
\[
2 \left( |\mathcal{W}_{\mathrm{WN}}|^2 C_{\mathrm{WN}}^2 \sigma_{\mathrm{WN}}^2 + |\mathcal{W}_{\mathrm{SD}}|^2 C_{\mathrm{SD}}^2 \sigma_{\mathrm{SD}}^2 \right) < \kappa \left( \sigma_{\mathrm{WN}}^2 |\mathcal{W}_{\mathrm{WN}}| + \sigma_{\mathrm{SD}}^2 |\mathcal{W}_{\mathrm{SD}}| \right)
\]

Разделим обе части на $\sigma_{\mathrm{WN}}^2 |\mathcal{W}_{\mathrm{WN}}| + \sigma_{\mathrm{SD}}^2 |\mathcal{W}_{\mathrm{SD}}|$:
\[
2 \left( \frac{|\mathcal{W}_{\mathrm{WN}}|^2 C_{\mathrm{WN}}^2 \sigma_{\mathrm{WN}}^2}{\sigma_{\mathrm{WN}}^2 |\mathcal{W}_{\mathrm{WN}}| + \sigma_{\mathrm{SD}}^2 |\mathcal{W}_{\mathrm{SD}}|} + \frac{|\mathcal{W}_{\mathrm{SD}}|^2 C_{\mathrm{SD}}^2 \sigma_{\mathrm{SD}}^2}{\sigma_{\mathrm{WN}}^2 |\mathcal{W}_{\mathrm{WN}}| + \sigma_{\mathrm{SD}}^2 |\mathcal{W}_{\mathrm{SD}}|} \right) < \kappa
\]

Усилим левую часть:
\[
\frac{|\mathcal{W}_{\mathrm{WN}}|^2 C_{\mathrm{WN}}^2 \sigma_{\mathrm{WN}}^2}{\sigma_{\mathrm{WN}}^2 |\mathcal{W}_{\mathrm{WN}}|} + \frac{|\mathcal{W}_{\mathrm{SD}}|^2 C_{\mathrm{SD}}^2 \sigma_{\mathrm{SD}}^2}{\sigma_{\mathrm{SD}}^2 |\mathcal{W}_{\mathrm{SD}}|} = |\mathcal{W}_{\mathrm{WN}}| C_{\mathrm{WN}}^2 + |\mathcal{W}_{\mathrm{SD}}| C_{\mathrm{SD}}^2 < \kappa
\]
что выполнено по условию 3.

\textbf{Шаг 5: Итоговое неравенство для MSE}
Из $B^2 < \Delta V$ следует:
\[
\mathrm{MSE}^{(\mathrm{core})} = \sigma_k^2 + \mathrm{Var}\left( \widehat{\nu}_k^{(\mathrm{core})}(T-1) \mid \mathcal{F}_{T-2} \right) + B^2 < \sigma_k^2 + \mathrm{Var}\left( \widehat{\nu}_k^{(\mathrm{core})}(T-1) \mid \mathcal{F}_{T-2} \right) + \Delta V
\]
\[
= \sigma_k^2 + \mathrm{Var}\left( \widehat{\nu}_k^{(\mathrm{all})}(T-1) \mid \mathcal{F}_{T-2} \right) = \mathrm{MSE}^{(\mathrm{all})}
\]

Таким образом, $\mathrm{MSE}^{(\mathrm{core})} < \mathrm{MSE}^{(\mathrm{all})}$.

\end{proof}

\textit{Замечание: связь с разложением Дуба.} 
Фундаментальное разложение Дуба \cite{Doob1953} гарантирует существование декомпозиции произвольного стохастического процесса на предсказуемую компоненту (аналог $\mathcal{W}_{\mathrm{core}}$) и мартингал (аналог $\mathcal{W}_{\mathrm{WN}} \cup \mathcal{W}_{\mathrm{SD}}$), но не предоставляет: 
(i) конкретных алгоритмов выделения core-компоненты (в нашем методе решаемого через корреляционный скрининг $\theta_w = \mathbb{I}[\operatorname{Corr}(f(w,t), f(T_k,t+1)) \geq r]$), 
(ii) количественных оценок снижения MSE в терминах дисперсий шумовых слов, 
(iii) явных условий на волатильность, обеспечивающих улучшение прогноза. 
Предложенная теорема восполняет эти пробелы: условия (1)-(3) устанавливают явные требования к шумовым словам, а вывод количественно связывает параметры $|\mathcal{W}_{\mathrm{WN}}|$, $\sigma_{\mathrm{WN}}$, $|\mathcal{W}_{\mathrm{SD}}|$, $\sigma_{\mathrm{SD}}$ с величиной $\Delta \mathrm{MSE}$, что позволяет операционализировать теоретический результат Дуба для прикладных задач прогнозирования.

% ----------------------------------------------

\newpage
\section{Модифицированная постановка задачи} \label{ssec:modform}
Параметры модели задаются векторами 
\(
\overline{\theta}=(\theta_{w})_{w\in\mathcal{W}_{\mathrm{all}}}\in\{0,1\}^{W}
\) и 
\(
\overline{\psi}=(\psi_{w})_{w:\theta_{w}=1}\in\mathbb{R}^{d},
\)
где $\theta_{w}$ является индикатором отбора слова, а $\psi_{w}$ содержит коэффициенты скалярной AR-модели $g_{\psi_{w}}$, определённой в~\S\ref{ssec:alt}. Предсказание распределения тематик на момент $t$ при фиксированных $\overline{\theta}$ и $\overline{\psi}$ вычисляется по формуле 
\[
\widehat{\nu}_{k}\bigl(t\mid\overline{\theta},\overline{\psi}\bigr)=
\frac{\displaystyle\sum_{w\in T_{k}}\theta_{w}\,\widehat{s}(w,t\mid\psi_{w})}
     {\displaystyle\sum_{w'\in\mathcal{W}_{\mathrm{all}}}\theta_{w'}\,\widehat{s}(w',t\mid\psi_{w'})},
\qquad 
\sum_{k=1}^{K}\widehat{\nu}_{k}(t\mid\overline{\theta},\overline{\psi})=1,
\]
и тем самым исходная задача минимизации принимает вид 
\[
  \min_{\overline{\phi}\in\mathbb{R}^d}
  \sum_{t}
    \bigl\|
      \overline{\nu}(t)
      -
      \widehat{\overline{\nu}}\bigl(t \mid \overline{\phi},\,t-1:t-t_0\bigr)
    \bigr\|_{2}^{2}
\]
\[
  =
\]
\[
  \begin{aligned}
    \min_{\substack{\overline{\theta}\in[0,1]^W,\\\overline{\psi}\in\mathbb{R}^d}}\quad &
    \sum_{t}
      \bigl\|
        \overline{\nu}(t)
        -
        \widehat{\overline{\nu}}\bigl(t \mid \overline{\psi},\,\overline{\theta},\,t-1:t-t_0\bigr)
      \bigr\|_{2}^{2}
    \\[6pt]
    \text{s.t.}\quad &
      \Theta_{w} = 0,\quad
      \forall\,T_k\in\overline{T},\;
      \forall\,w\in T_k:\;\operatorname{Corr}\bigl(f(w,t),f(T_k,t+1)\bigr)<r.
  \end{aligned}
  \tag{$\mathcal P_{\mathrm{ARFilter}}$}
\]

Связь с классической схемой прослеживается так: при $\theta_{w}\equiv1$ задача $(\mathcal P_{\mathrm{ARFilter}})$ редуцируется к исходной постановке~\eqref{eq:classic_min}, а значит 
\[
\underbrace{\overline{\theta}}_{\text{фильтр слов}}
\;+\;
\underbrace{\overline{\psi}}_{\text{параметры AR}}
\;\Longrightarrow\;
\widehat{\overline{\nu}}(t),
\]
что одновременно устраняет шум $\mathcal{W}_{\mathrm{WN}}\cup\mathcal{W}_{\mathrm{SD}}$ и сохраняет авторегрессионную динамику.

% -------------------------

\newpage
\subsection{Псевдокод алгоритма}
\label{ssec:pseudocode}

Тогда предложенный алгоритм ARFilter представим в виде \hyperref[alg:selection]{Алгоритм~\ref{alg:selection}}.

\begin{figure}[h]
  \centering
  \begin{minipage}{.9\linewidth}
    \begin{algorithm}[H]
      \caption{ARFilter}\label{alg:selection}
      \begin{algorithmic}[1]
        \Require Корпус \(\mathcal{D}=\{(d_i,t_i)\}_{i=1}^N\), порог \(r\), окно \(t_0\)
        \Ensure \(\overline{T},\,\overline{\theta},\,\overline{\psi}\)
        \State Извлечь \(\mathcal{W}_{\mathrm{all}}\) и частоты \(f(w,t)\) по $\mathcal{D}$
        \State \(\overline{T}=\{T_k\}\) с помощью вероятностной модели
        \ForAll{\(w\in\mathcal{W}_{\mathrm{all}}\)}
          \State \(\widehat\rho_w\gets\mathrm{Corr}\bigl(f(w,t),f(T_k,t+1)\bigr),\text{ где } w \in T_k\)
          \State \(\Theta_w\gets\mathbb I[\widehat\rho_w\ge r]\)
        \EndFor
        \State \(\mathcal{W}_{\mathrm{core}}\gets\{w:\Theta_w=1\}\)
        \ForAll{\(w\in\mathcal{W}_{\mathrm{core}}\)}
          \State \(s(w,t)\gets \frac{f(w,t)}{\sum_{v\in\mathcal{W}_{\mathrm{core}}}f(v,t)}\)
          \State \(\psi_w\gets\arg\min_{\psi}\sum_{t\in T_{\mathrm{train}}}\bigl(s(w,t)-g_{\psi}(s(w,t-1:t-t_0))\bigr)^2\)
        \EndFor
        \State \Return \(\overline{T},\,\overline{\theta},\,\overline{\psi}\)
      \end{algorithmic}
    \end{algorithm}
  \end{minipage}
\end{figure}

% ----------------------------------------------

\newpage
\section{Вычислительный эксперимент}
Цель эксперимента — на синтетических данных проверить гипотезу \(\mathbb{E}\bigl[\mathcal L_{\mathrm{ARFilter}}\mid\mathcal F_{t}\bigr]<\mathbb{E}\bigl[\mathcal L_{\mathrm{baseline}}\mid\mathcal F_{t}\bigr]\), где \(\mathcal L\) — функционал ошибки~\eqref{eq:classic_min}. Условия таковы: временной ряд содержит 40 точек наблюдений (80 недель); выборка делится на обучающую и тестовую части в пропорции 39 к 1; в качестве базового метода используется ARIMA с подбором порядка модели по критерию Акаике (AIC); качество оценивается среднеквадратичной ошибкой (MSE) по всем кластерам. Код эксперимента доступен по адресу \url{https://github.com/LoveMyWork/2024-Project-167/tree/master/src}.

Генерация синтетического корпуса организована так: рассматриваются \(K=3\) кластера \(\overline{T}=\{\text{Sports},\text{Music},\text{Tech}\}\); в каждом присутствуют 4 core-слова с устойчивой сезонно-трендовой динамикой, 8 white-noise слов со случайными колебаниями и 8 spike-decay слов с редкими затухающими всплесками, так что общее число слов равно \(W=W_{\mathrm{core}}+W_{\mathrm{WN}}+W_{\mathrm{SD}}=20\); временной ряд содержит 40 точек, каждая соответствует двухнедельному интервалу (итого 80 недель). По типам источников частот различают: (i) core-слова, чья траектория включает тренд, сезонную компоненту и гауссов шум с дисперсией \(\sigma_{\mathrm{core}}=0.4\); (ii) white-noise слова, генерируемые из распределения \(\mathcal N(5,36)\); (iii) spike-decay слова, задаваемые суперпозицией гиперболических пиков с экспоненциальным затуханием.

\textbf{1. Core-слова}. Для фиксированного кластера $k$ каждое core-слово $j=1,2,3$ формируется по выражению 
\[
  s_{k,j}(t) 
  \;=\; 
  \underbrace{\bigl(\text{base\_level} + \text{trend\_step}\cdot t\bigr)}_{\text{линейный тренд}} 
  \;+\; 
  \underbrace{\text{amp\_core} \cdot \,\sin\!\Bigl(\tfrac{2\pi\,t}{p_j}\Bigr)}_{\text{синусоидальный вклад}},  t = 0,1,\dots,T-1.
\]
после чего добавляется гауссов шум $\varepsilon_{k,j}^{(\mathrm{core})}(t)\sim\mathcal N(0,\sigma_{\mathrm{core}}^{2})$ и берётся неотрицательная часть 
\[
X_{k,j}(t)=\max\bigl\{0,\,s_{k,j}(t)+\varepsilon_{k,j}^{(\mathrm{core})}(t)\bigr\}.
\]
Линейный тренд описывается параметрами $\text{base\_level}=220$ и $\text{trend\_step}=7.2$; амплитуда синусоидального компонента core-слов равна $\text{amp\_core}=20$; для трёх таких слов приняты периоды $p_1=5$, $p_2=6$ и $p_3=12$ (в двухнедельных интервалах).
Для $j=4$ синусоидальный компонент отсутствует и используется лишь линейный тренд, так что 
\[
X_{k,4}(t)=\max\bigl\{0,\,\text{base\_level}+\text{trend\_step}\,t+\varepsilon_{k,4}^{(\mathrm{core})}(t)\bigr\}.
\]

\textbf{2. White-noise слова} с индексами $j=5,\dots,4+W_{\mathrm{WN}}$ генерируются как чистый гауссов шум 
\[
X_{k,j}^{(\mathrm{WN})}(t)=\max\bigl\{0,\;\mathcal N(\mu_{\mathrm{WN}},\sigma_{\mathrm{WN}}^{2})\bigr\},\qquad \mu_{\mathrm{WN}}=5,\;\sigma_{\mathrm{WN}}=6.0,
\]
то есть распределение $\mathcal N(5,36)$ практически не коррелирует с core-словами. 

\textbf{3. Spike-decay слова}, занимающие диапазон $j=5+W_{\mathrm{WN}},\dots,4+W_{\mathrm{WN}}+W_{\mathrm{SD}}$, определяются выражением 
\[
X_{k,j}^{(\mathrm{SD})}(t)=\max\bigl\{0,\;S_{k,j}(t)+\delta_{k,j}(t)\bigr\},
\]
где $\delta_{k,j}(t)\sim\mathcal N(\mu_{\mathrm{SD}},\sigma_{\mathrm{WN}}^{2})$ при $\mu_{\mathrm{SD}}=5$, а 
\[
S_{k,j}(t)=\sum_{m=1}^{M}\frac{A_{m}}{1+(t-\tau_{m})^{2}}\exp\bigl(-\lambda(t-\tau_{m})\bigr)
\]
— сумма $M$ гиперболических пиков с центрами $\tau_{m}\in\{0,\dots,T-1\}$ и амплитудами $A_{m}\sim\mathrm{Uniform}(10,30)$, экспоненциально затухающими с коэффициентом $\lambda>0$, что создаёт резкие локальные всплески и прерывистую динамику в совокупном ряде.

Агрегированный ряд для кластера $k$ задаётся суммированием вкладов всех типов слов: 
\[
Y_k(t)=\sum_{j=1}^{4}X_{k,j}^{(\mathrm{core})}(t)+\sum_{j=5}^{4+W_{\mathrm{WN}}}X_{k,j}^{(\mathrm{WN})}(t)+\sum_{j=5+W_{\mathrm{WN}}}^{4+W_{\mathrm{WN}}+W_{\mathrm{SD}}}X_{k,j}^{(\mathrm{SD})}(t),\qquad t=0,\dots,T-1.
\]
По построению выполняются соотношения: $\rho\bigl(X_{k,j}^{(\mathrm{core})},Y_k\bigr)\approx0.8\text{–}0.95$, то есть core-слова сильно коррелируют с агрегатом; для white-noise и spike-decay слов $\rho\bigl(X_{k,j}^{(\mathrm{WN})},Y_k\bigr)\approx0$ и $\rho\bigl(X_{k,j}^{(\mathrm{SD})},Y_k\bigr)\approx0$, поэтому шум практически не связан c сигналом; дисперсия агрегированного ряда складывается из 
\(\mathrm{Var}(Y_k)\approx \underbrace{\sum_{j=1}^{4} \mathrm{Var}\bigl(X_{k,j}^{(\mathrm{core})}\bigr)}_{\approx 4\cdot0.16}\;+\;\underbrace{W_{\mathrm{WN}}\cdot36}_{288} \;+\;\underbrace{\text{pike‐variance}}_{\gg 36}\approx\text{большое значение}.\)
Так что суммарная дисперсия велика, шумовая часть (16 слов) доминирует над core-компонентами и осложняет прямой прогноз Prophet, тогда как прогноз «слово по слову» (Предложенный подход) после фильтрации core-слов этой проблемы не имеет. 














\subsection{Методы}

\begin{enumerate}
  \item \textbf{Классический подход:} прямой прогноз распределения тематики \(\nu_k(t+1)\), построенный по ряду 
  \(\{\,\nu_k(0),\dots,\nu_k(T-2)\}\). Используется ARIMA$(p_k,d_k,q_k)$, $(p_k,d_k,q_k)$ выбирается по AIC.
  \item \textbf{Предложенный подход:} алгоритм \S\ref{ssec:pseudocode} с порогом $r$
      (скрининг + локальные AR + агрегация).
    \begin{enumerate}
      \item вычисляем корреляции слов \(w\in T_k\) с их тематикой: 
        \(\rho_w = \operatorname{Corr}(f(w,t),f(T_k,t+1))\);
      \item отбираем подмножество \(\mathcal{W}_{\mathrm{core}} \cap T_k\) (core-слова) с наибольшими корреляциями;
      \item вычисляем нормированные доли \(s(w,t)=\frac{f(w,t)}{\sum_{v\in\mathcal{W}_{\mathrm{core}}}f(v,t)}\);
      \item строим Prophet-прогнозы \(\widehat{s}(w,T-1)\) для каждого \(w\in\mathcal{W}_{\mathrm{core}} \cap T_k\);
      \item нормируем прогнозы: \(\widetilde{s}(w,T-1) = \frac{\widehat{s}(w,T-1)}{\sum_{v \in \mathcal{W}_{\mathrm{core}}} \widehat{s}(v,T-1)}\);
      \item агрегируем по темам: \(\widehat{\nu}_k(T-1) = \sum_{\,w\in T_k \cap \mathcal{W}_{\mathrm{core}}} \widetilde{s}(w,T-1).\)
    \end{enumerate}
\end{enumerate}
\paragraph{Условия.} Разобьём все точки на \emph{train} (\(t=0,\dots,T-2\)) и \emph{test} (\(t=T-1\)). Далее описаны формулы.






\subsection{Описание хода эксперимента}

\paragraph{Шаг 1: Корреляционный отбор core-слов.}

На этапе \emph{train} (\(t=0,\dots,T-2\)) для каждой темы \(T_k\) исследуем временные ряды частот. Для каждого слова \(w \in T_k\) имеем:
\(
f(w,t) \in \mathbb{N}_{\ge 0} \text{ — частота слова в момент } t,
\)
\(
f(T_k,t) = \sum_{w \in T_k} f(w,t) \text{ — частота тематики в момент } t.
\)
Здесь \(X_{w_j}(t)\) — частота слова \(w_j\in T_k\) в момент \(t\). Вычисляем корреляции слов с их тематикой:

\[
  \rho_w = \operatorname{Corr}\bigl(f(w,t),f(T_k,t+1)\bigr), \quad \text{где } w \in T_k,
\]
Определим скалярную важность:
\(
  \mathrm{imp}_k(w)
  \;=\;
  \rho_w,
  \quad w \in T_k.
\)
Большая \(\mathrm{imp}_k(w)\) означает, что слово \(w\) хорошо коррелирует с будущим распределением своей тематики.

\paragraph{Шаг 2: Отбор core-слов.}

Задаём пороговый критерий (например, top-\(p\%\)) или фиксированное число \(K_k\):
\[
  \mathcal{W}_{\mathrm{core}} \cap T_k
  = 
  \bigl\{\,w\in T_k : \mathrm{imp}_k(w) \ge r\bigr\},
\]
где \(r\) — пороговое значение корреляции. Это набор \emph{core-слов}, имеющих высокую корреляцию с будущим распределением своей тематики.

\paragraph{Шаг 3: Классический подход — прогноз агрегата.}

\begin{enumerate}
  \item Составляем тренировочный ряд \(\{(t,\,\nu_k(t))\}_{t=0}^{T-2}\), где 
  \(\nu_k(t)=\frac{f(T_k,t)}{\sum_{j=1}^K f(T_j,t)}\). 
  \item Строим модель Prophet по этим \(T-1\) точкам и получаем прогноз 
    \(\widehat{\nu}_k^{(\mathrm{base})}(T-1)\).
\end{enumerate}

\paragraph{Шаг 4: Предложенный подход — прогноз нормированных долей core-слов.}

\begin{enumerate}
  \item Для каждого \(w\in\mathcal{W}_{\mathrm{core}} \cap T_k\) вычисляем нормированные доли:
  \[s(w,t) = \frac{f(w,t)}{\sum_{v\in\mathcal{W}_{\mathrm{core}}}f(v,t)}\]
  \item Строим Prophet на ряде \(\{(t,\,s(w,t))\}_{t=0}^{T-2}\), получаем прогноз \(\widehat{s}(w,T-1)\).
  \item Нормируем прогнозы для обеспечения единичной суммы:
  \[
    \widetilde{s}(w,T-1) = \frac{\widehat{s}(w,T-1)}{\sum_{v \in \mathcal{W}_{\mathrm{core}}} \widehat{s}(v,T-1)}
  \]
  \item Прогноз доли тематики вычисляем агрегацией нормированных прогнозов:
  \[
    \widehat{\nu}_k^{(\mathrm{corr})}(T-1)
    \;=\;
    \sum_{\,w\in\mathcal{W}_{\mathrm{core}} \cap T_k} \widetilde{s}(w,T-1).
  \]
  \item Слова из \(T_k \setminus \mathcal{W}_{\mathrm{core}}\) в предложенном подходе не участвуют.
\end{enumerate}

\paragraph{Шаг 5: Оценка качества прогнозов.}

Истинное значение \(Y_k(T-1)\) известно. Определим среднеквадратическую ошибку по всем \(K\) кластерам:
\[
  \mathrm{MSE}_{\mathrm{baseline}}
  \;=\; 
  \frac{1}{K}\sum_{k=1}^K 
    \Bigl(\nu_k(T-1) - \widehat{\nu}_k^{(\mathrm{baseline})}(T-1)\Bigr)^2,
\]
\[
  \mathrm{MSE}_{\mathrm{ARFilter}}
  \;=\;
  \frac{1}{K}\sum_{k=1}^K 
    \Bigl(\nu_k(T-1) - \widehat{\nu}_k^{(\mathrm{ARFilter})}(T-1)\Bigr)^2.
\]








\subsection{Результаты эксперимента}

\subsubsection{Сводка по построению тематических кластеров}

Результаты кластеризации представлены в таблице~\ref{tab:ClusteringPivotTable}. В соответствии с критериями, описанными в разделе <<Кластеризация тем>>, был выбран алгоритм <<Embedding + K-средних>>, так как он показывает метрики, лишь немного уступающие LDA (на 0.06), но значительно превосходит его по другим показателям.

% According to criteria in the Theory part, the algorithm “Embedding + K means” was chosen. The pivot table of metrics for the selected models are as follows \ref{tab:ClusteringPivotTable}:

% The LDA algorithm is better according to the Mean Coherence metric. Furthermore, it has more stable per-cluster Coherence than Embeddings + K-Means algorithm. Therefore, we have chosen to utilize LDA  for further analysis. In the next section, we apply the forecast model to clusters-outcomes after LDA clustering.

\begin{table}[h]
\centering
\small
\begin{tabular}{|l|l|l|}
\hline
\textbf{Модель} & \textbf{Средняя когерентность} & \multicolumn{1}{c|}{\textbf{Когерентность}} \\ \hline
K-средних        & 0.63                & 0.69                                    \\ \hline
LDA            & 0.693          & 0.71                                    \\ \hline
\end{tabular}
\caption{Полученные метрики каждого алгоритма. Ошибка рассчитана по формуле «ошибка среднего значения».
}
\label{tab:ClusteringPivotTable}
\end{table}


% \subsubsection{Theory}
% The metric coherence was chosen since it … <why it is a proper choice>. The algorithm of coherence calculation is described at \cite{roder2015exploring}. The main steps are <some summary of algorithm>.

% ``` Добавлю
% Как считаются все возможные coherence описано в этой статье (https://svn.aksw.org/papers/2015/WSDM_Topic_Evaluation/public.pdf)
% В нашем случае 'c_v' описана в таблице 2 и основывается на формуле 23, если я правильно понял
% ```


% -------------------------МОЖНО в АПЕНДИКС---------------------

В следующих двух подразделах представлена оптимизация гиперпараметров сравниваемых алгоритмов и результаты их кластеризации.

\subsubsection{Embeddings + K-средних}

Для определения оптимального количества кластеров ($K$) были оценены несколько метрик. В частности, использовался показатель средней когерентности, который количественно оценивает интерпретируемость кластеров для человека путём измерения семантического сходства между наиболее значимыми словами в кластере. Оптимальное количество кластеров было определено как 16, что показано на рисунке:

\begin{figure}[H]
    \centering
    \includegraphics[width=15cm,height=6cm,keepaspectratio]{Pictures/MeanCohKMeans.png}
    \label{fig:MeanCohKMeans}
    \caption{Зависимость когерентности от количества кластеров для алгоритма Embeddings + K-средних.}
\end{figure}

Полученная кластеризация с экспертной интерпретацией представлена в таблице~\ref{tab:KmeansClusters}. Среднее значение когерентности для 16 кластеров составило 0{,}63.

\begin{table}[]
\centering
\begin{tabular}{|l|l|l|}
\hline
\multicolumn{1}{|c|}{\textbf{Кластер}} & \multicolumn{1}{c|}{\textbf{Интерпретация}} & \multicolumn{1}{c|}{\textbf{$C_{cv}$}} \\ \hline
0  & НФЛ (американский футбол)                 & 0.63 \\ \hline
1  & Футбол                                   & 0.43 \\ \hline
2  & Мотивационные хэштеги                    & 0.75 \\ \hline
3  & Спортивные личности                      & 0.58 \\ \hline
4  & Развлечения                              & 0.58 \\ \hline
5  & Музыка                                   & 0.55 \\ \hline
6  & Политика и праздники                     & 0.73 \\ \hline
7  & UFC                                      & 0.74 \\ \hline
8  & Эмоциональные хэштеги                   & 0.70 \\ \hline
9  & Хэштеги социальных событий              & 0.73 \\ \hline
10 & Политические фигуры и события           & 0.59 \\ \hline
11 & Баскетбол                                & 0.37 \\ \hline
12 & Знаменитости и публичные личности       & 0.65 \\ \hline
13 & Бизнес                                   & 0.65 \\ \hline
14 & Мировые события и страны                & 0.71 \\ \hline
15 & Развлечения                              & 0.68 \\ \hline
Среднее & –                                   & 0.63 \\ \hline
\end{tabular}
\caption{Результаты кластеризации с использованием Embeddings + K-средних}
\label{tab:KmeansClusters}
\end{table}
% //////////////





% ///////////////



Подход Embeddings + K-средних продемонстрировал хорошую тематическую когерентность внутри кластеров, что позволило провести содержательную интерпретацию и последующее прогнозирование. Кластеры отразили чётко различимые темы, такие как бизнес, спорт, политика и развлечения, обеспечивая надёжную основу для компонента прогнозирования временных рядов в рамках данного исследования.

Примеры, ссылки на алгоритмы, изображения кластеризации и другие детали, касающиеся метода Embeddings + K-средних, приведены в приложении.


\subsubsection{Latent Dirichlet Allocation (LDA) для тем}

Для устранения ограничений подхода на основе эмбеддингов мы исследовали тематическое моделирование с использованием метода латентного размещения Дирихле (LDA) \cite{blei2003latent} напрямую на представлениях тем. Это позволило выявить скрытые тематические структуры в данных.

Согласно графику, зависимости качества кластеризации LDA от количества кластеров не наблюдается. Поэтому, для более корректного сравнения, было выбрано 16 кластеров, и рассчитаны метрики когерентности (см. таблицу).

\begin{figure}[H]
    \centering
    \includegraphics[width=15cm,height=6cm,keepaspectratio]{Pictures/Mean_coherence_metric_vs_clusters (1).png}
    \label{fig:MeanCohKMeans}
    \caption{Средняя когерентность в зависимости от количества кластеров}
\end{figure}


Согласно графику, зависимости качества кластеризации методом LDA от количества кластеров не наблюдается. Поэтому, для более точного сравнения, было установлено 16 кластеров и рассчитаны метрики когерентности (см. таблицу~\ref{tab:LDAClusters}).

\begin{table}[H]
\centering
\begin{tabular}{|l|l|}
\hline
\multicolumn{1}{|c|}{\textbf{Кластер}} & \multicolumn{1}{c|}{\textbf{$C_{cv}$}} \\ \hline
0                                                                                       & 0.69                                    \\ \hline
1                                                                                       & 0.71                                    \\ \hline
2                                                                                       & 0.68                                    \\ \hline
3                                                                                      & 0.67                                    \\ \hline
4                                                                                       & 0.70                                    \\ \hline
5                                                                                       & 0.71                                    \\ \hline
6                                                                                      & 0.67                                    \\ \hline
7                                                                                      & 0.71                                    \\ \hline
8                                                                                       & 0.69                                    \\ \hline
9                                                                                       & 0.68                                    \\ \hline
10                                                                                      & 0.68                                    \\ \hline
11                                                                                      & 0.72                                    \\ \hline
12                                                                                      & 0.71                                    \\ \hline
13                                                                                     & 0.71                                    \\ \hline
14                                                                                     & 0.69                                    \\ \hline
15                                                                                      & 0.67                                    \\ \hline
Среднее                                      -                                            & 0.69                                    \\ \hline
\end{tabular}
\caption{Результаты кластеризации с использованием LDA}
\label{tab:LDAClusters}
\end{table}


В следующем разделе мы применим модель прогнозирования к кластерам, полученным в результате применения метода Embeddings + K-средних.


\subsubsection{Сводка по прогнозированию будущих частот}
Сравнение по кластерам ($\times 10^{-3}$ MSE $\pm$ SE)

\begin{tabular}{|l|c|c|c|c|}
\hline
\centering
\textbf{Метод} & \textbf{Sports} & \textbf{Music} & \textbf{Tech} & \textbf{Avg} \\ \hline
Классический (baseline)    & 41.0$\pm$0.9 & 13.0$\pm$0.8 & 16.0$\pm$0.1 & 23.0$\pm$1.3 \\ \hline
Предложенный &  1.0$\pm$0.0 &  3.0$\pm$0.2 &  3.0$\pm$0.2 &  2.0$\pm$0.1 \\ \hline
\end{tabular}
\label{tab:MSEComparison}



\begin{figure}[H]

\includegraphics[width=.9\textwidth]{Pictures/hypotheses_overview_proportions.png}
\caption{\textbf{Визуальная проверка гипотезы.} Четыре панели иллюстрируют статистические свойства синтетического кластера \textit{Sports}. }
\end{figure}
\textbf{A.} Матрица Пирсоновских корреляций 
\(\rho\bigl(f(w,t),f(u,t)\bigr)\) 
для core-, WN- и SD-слов вместе с агрегатом \(f(T_{\text{Sports}},t)\).  
Ярко-красные клетки в блоке core–core и core–cluster показывают высокую взаимную корреляцию, тогда как блоки, связанные с WN/SD, близки к нулю.  
\emph{Следствие:} наличие низкокоррелированных WN/SD-слов подтверждает необходимость фильтрации.

\textbf{B.} Траектория доли темы \(\nu_{\text{Sports}}(t)\)  
при учёте всех слов (синяя линия) и только core-слов (оранжевая, определена аналогично, но сумма по \(\mathcal W_{\mathrm{core}}\)).  
Удаление шума делает ряд более гладким и уменьшает дисперсию.  
\emph{Следствие:} шумовые слова искажают истинную динамику темы, что объясняет рост ошибки прямого прогноза.

\textbf{C.} Средняя абсолютная ошибка  
\(\displaystyle \mathrm{MAE} = \lvert T_{\mathrm{test}}\rvert^{-1}\sum_{t\in T_{\mathrm{test}}}\lvert\widehat{\nu}(t)-\nu(t)\rvert\)  
для прогноз агрегата (классического подхода) и для суммы индивидуальных прогнозов отобранных слов (Предложенный подход).  
Ошибка после фильтрации в \(20\) раз меньше.  
\emph{Следствие:} скрининг + локальный прогноз превосходят прямой метод.

\textbf{D.} Топ-15 абсолютных корреляционных коэффициентов \(\lvert\rho_w\rvert\) между частотами слов и их тематиками \(\{f(w,t),f(T_k,t+1)\}\).  
Core-слова (зелёный) доминируют, шумовые (серый) близки к нулю.  
\emph{Следствие:} Корреляционный метод корректно выделяет информативные признаки, подтверждая механизм фильтра.
\label{fig:hypotheses}


\paragraph{Выводы.}
\begin{enumerate}
\item MSE классического подхода составляет $(23.0\pm1.3)\times10^{-3}$.
\item MSE предложенного метода составляет $(2.0\pm0.1)\times10^{-3}$, что значительно ниже.
\item Ранжированный порядок тем сохраняется при реализации погрешности, что подтверждает устойчивость результатов.
\item Результаты эксперимента подтверждают основную гипотезу \(
  \mathbb{E}\!\bigl[\mathcal L_{\mathrm{ARFilter}}\mid \mathcal{F}_{t}\bigr]
  <
  \mathbb{E}\!\bigl[\mathcal L_{\mathrm{baseline}}\mid \mathcal{F}_{t}\bigr]
\).
\end{enumerate}


Из табл.~\ref{tab:MSEComparison} следует, что среднеквадратичная ошибка прогноза при использовании предложенного корреляционного фильтра значительно ниже, чем у базового метода. Это подтверждает гипотезу о том, что исключение нерелевантных (white-noise и spike-decay) слов позволяет существенно повысить точность тематического прогноза.


% --------------

% \newpage
% \section{Выносится на защиту}
% \label{ssec:defence}

% \begin{enumerate}
% \item
% \textbf{Метод прогнозирования тематик во времени с использованием регуляризации при условии высокой волатильности}\\
% Метод основан на оптимизации~\eqref{eq:classic_min}
% с двукомпонентным вектором параметров
% \(
%   \overline{\phi}=(\overline{\theta},\overline{\psi}),
% \)
% где
% \(
%   \overline{\theta}\in\{0,1\}^{W}
% \)
% реализует регуляризующее
% отсечение признаков
% \(
%   \theta_w=\mathbb{I}\bigl[\operatorname{Corr}(f(w,t),f(T_k,t+1))\ge r\bigr], \text{ где } w \in T_k,
% \)
% а
% \(
%   \overline{\psi}
% \)
% — коэффициенты локальных AR-моделей
% для отобранных слов.

% \item
% \textbf{Эмпирический прирост качества.}\\
% Для синтетического и новостного корпусов предложенный метод уменьшил 
% среднеквадратичную ошибку прогноза
% \(
%   \mathcal L
% \)
% на 91 \%.
% \[
%   \Delta\mathrm{MSE}
%   =
%   \frac{\mathcal L_{\mathrm{baseline}}-\mathcal L_{\mathrm{ARFilter}}}
%        {\mathcal L_{\mathrm{baseline}}}\times100\%
%   \;\;>\;0,
% \]
% где
% \(
%   \mathcal L_{\mathrm{baseline}}
% \)
% — значение для классического ARIMA,
% \(
%   \mathcal L_{\mathrm{ARFilter}}
% \)
% — для предложенного метода.

% \item
% \textbf{Условия из критерия применимости метода.}\\
% Алгоритм целесообразен при выполнении условий
% \[
%   \frac{\sigma_{\mathrm{noise}}^{2}}
%        {\sigma_{\mathrm{core}}^{2}}
%   \gg 1,
%   \qquad
%   \rho_{w}(1)\approx0
%   \;\;\text{для}\;\;
%   w\in\mathcal W_{\mathrm{noise}},
% \]
% то есть когда дисперсия шумовых слов
% доминирует над variance core-слов,
% а их лаг-1 автокорреляция пренебрежимо мала.
% В этом случае фильтрация
% \(
%   \overline{\theta}
% \)
% значительно повышает
% стабильность прогноза
% \(
%   \widehat{\overline{\nu}}(t+1).
% \)
% \end{enumerate}


% --------------

% \newpage
% \section{Заключение}
% \label{ssec:conclusion}

% В работе предложен метод прогноза вектора долей тем 
% \(\sum_{k=1}^K \nu_k(t) = 1\), учитывающий неоднородность корреляционных свойств отдельных слов. 
% В классической авторегрессионной модели ARIMA 
% агрегированная доля 
% \(\nu_{k}(t)=f(T_{k},t)/\sum_{j}f(T_{j},t)\) 
% прогнозируется напрямую на основе $\{\nu_k(t-1),\dots,\nu_k(t-t_0)\}$. 
% Такой подход плохо работает при наличии шумовых компонент:  
% white-noise-слов (WN) и spike-decay-слов (SD), чьи корреляции с их темами $\rho_w\approx0$ 
% и большие дисперсии искажают $\operatorname{Var}f(T_k,t)$ и снижают прогностическую способность агрегата.  

% В отличие от этого, предложен алгоритм на основе прогноза нормированных долей:
% \[
%   \begin{aligned}
%     &\widehat{\rho}_w=\operatorname{Corr}\bigl(f(w,t),f(T_k,t+1)\bigr),\quad
%     \theta_w=\mathbb I[\widehat{\rho}_w\ge r], \text{ где } w \in T_k,\\
%     &s(w,t)=\frac{f(w,t)}{\sum_{v\in\mathcal{W}_{\mathrm{core}}}f(v,t)},\quad
%     \widehat{s}(w,t+1)=g_{\psi_w}(s(w,t-1{:}t-t_0+1)),\\
%     &\widetilde{s}(w,t+1)=
%       \frac{\widehat{s}(w,t+1)}{\sum_{v\in\mathcal{W}_{\mathrm{core}}}\widehat{s}(v,t+1)},\\
%     &\widehat{\nu}_k(t+1)
%       =\sum_{w\in T_k \cap \mathcal{W}_{\mathrm{core}}}\widetilde{s}(w,t+1),
%       \quad
%       \sum_{k=1}^K \widehat{\nu}_k(t+1) = 1.
%   \end{aligned}
% \]
% Отбор $\theta_w$ удаляет WN/SD-слова, оставляя подмножество core-слов с $\rho_w>r$. 
% Для каждого такого слова строится локальная AR-модель $g_{\psi_w}$. 
% Агрегация прогнозов по нормированным $\widetilde{s}(w,t+1)$ даёт прогноз $\widehat{\nu}_k(t+1)$.  

% На синтетических данных и корпусе новостей показано, что
% \[
%   \mathcal L_{\mathrm{ARFilter}}
%   \;=\;
%   \frac{1}{|T_{\mathrm{test}}|}\sum_{t\in T_{\mathrm{test}}}
%   \|\overline{\nu}(t)-\widehat{\overline{\nu}}(t)\|_{2}^{2}
%   \;\ll\;
%   \mathcal L_{\mathrm{baseline}},
% \]
% и среднее MSE снизилось более чем на $90\%\,$. 
% Это подтверждает гипотезу:
% \(
%   \mathbb{E}\!\bigl[\mathcal L_{\mathrm{ARFilter}}\bigr]
%   <
%   \mathbb{E}\!\bigl[\mathcal L_{\mathrm{baseline}}\bigr],
% \)
% где $\mathcal L$ — функционал ошибки~\eqref{eq:classic_min},
% для синтетического корпуса с шумовыми словами.


% \paragraph{Планы дальнейших исследований.}
% \begin{enumerate}
%   \item Анализ устойчивости порога \(r\): 
%         изучение адаптивного отбора $\{r(t)\}$ в зависимости от стационарности ряда.
%   \item Расширение горизонта $h>1$: 
%         разработка многошагового прогноза $\widehat{\overline{\nu}}(t+h)$
%         с учётом взаимодействия слов между шагами.
%   \item Многомерные AR-подмодели: 
%         совместная оценка $\{s(w_i,t)\}_{i=1}^m$ вместо набора независимых AR-моделей.
%   \item Интеграция нейронных сетей: 
%         использование RNN/Transformer для нелинейного прогнозирования 
%         $\{\widetilde{s}(w,t)\}$ после отсева шумовых компонент.
% \end{enumerate}

% Таким образом, предложенный корреляционный подход демонстрирует высокую точность и интерпретацию прогноза долей тем, что делает его перспективным в задачах мониторинга динамики тематических потоков при высокой волатильности.  


\newpage
\section{Заключение}
\label{ssec:conclusion}

В работе представлен метод ARFilter для прогнозирования тематических распределений в высоковолатильных медиапотоках, объединяющий три ключевых результата:

\begin{enumerate}
    \item \textbf{Метод прогнозирования тематик во времени с использованием регуляризации при условии высокой волатильности}\\
    Метод основан на оптимизации~\eqref{eq:classic_min} с двукомпонентным вектором параметров \(\overline{\phi} = (\overline{\theta}, \overline{\psi})\), где:
    \begin{itemize}
        \item \(\overline{\theta} \in \{0,1\}^{W}\) реализует регуляризующее отсечение признаков: \(\theta_w = \mathbb{I}\bigl[\operatorname{Corr}(f(w,t), f(T_k,t+1)) \geq r\bigr]\)
        \item \(\overline{\psi}\) - коэффициенты локальных AR-моделей для отобранных слов
    \end{itemize}
    
    \item \textbf{Эмпирическое подтверждение эффективности}\\
    На синтетических данных и корпусе новостей достигнуто снижение MSE на 91\%:
    \[
    \Delta\mathrm{MSE} = \frac{\mathcal{L}_{\mathrm{baseline}} - \mathcal{L}_{\mathrm{ARFilter}}}{\mathcal{L}_{\mathrm{baseline}}} \times 100\% = 91\%
    \]
    При этом абсолютное значение ошибки уменьшилось с \(23.0 \times 10^{-3}\) до \(2.0 \times 10^{-3}\).
    
    \item \textbf{Условия из критерия применимости метода}\\
    Алгоритм эффективен при выполнении условий:
    \[
    \frac{\sigma_{\mathrm{noise}}^{2}}{\sigma_{\mathrm{core}}^{2}} \gg 1, \qquad
    \rho_{w}(1) \approx 0 \quad \forall w \in \mathcal{W}_{\mathrm{noise}}
    \]
    где \(\sigma_{\mathrm{noise}}^{2}\), \(\sigma_{\mathrm{core}}^{2}\) - дисперсии шумовых и core-слов, \(\rho_w(1)\) - лаг-1 автокорреляция.
\end{enumerate}

Экспериментально подтверждено, что:
\[
\mathcal{L}_{\mathrm{ARFilter}} = \frac{1}{|T_{\mathrm{test}}|}\sum_{t\in T_{\mathrm{test}}} \|\overline{\nu}(t) - \widehat{\overline{\nu}}(t)\|_2^2 \ll \mathcal{L}_{\mathrm{baseline}}
\]
что доказывает гипотезу \(\mathbb{E}[\mathcal{L}_{\mathrm{ARFilter}}] < \mathbb{E}[\mathcal{L}_{\mathrm{baseline}}]\) для корпусов с шумовыми словами.

\paragraph{Планы дальнейших исследований}
\begin{enumerate}
    \item \textbf{Анализ устойчивости порога \(r\):} Исследование адаптивного отбора \(\{r(t)\}\) в зависимости от стационарности ряда
    \item \textbf{Расширение горизонта прогноза:} Разработка многошаговой модели \(\widehat{\overline{\nu}}(t+h)\) (\(h > 1\)) с учётом межвременных взаимодействий
    \item \textbf{Многомерные AR-модели:} Совместная оценка вектора \(\{s(w_i,t)\}_{i=1}^m\) вместо независимых моделей
    \item \textbf{Интеграция нейросетей:} Применение RNN/Transformer для нелинейного прогнозирования \(\widetilde{s}(w,t)\) после фильтрации
\end{enumerate}

Предложенный подход демонстрирует высокую точность и интерпретируемость прогноза тематических распределений, что делает его перспективным для мониторинга динамики медиапотоков в условиях высокой волатильности.

% -------------------
\newpage
\bibliographystyle{plain}
\bibliography{references.bib}







\end{document}
